\documentclass[11pt]{article}
\usepackage[margin=1.1in]{geometry}
\usepackage{amsmath,amssymb,amsthm}
\usepackage{booktabs}
\usepackage{rossmeta}
\usepackage{doi}

\newtheorem{theorem}{Theorem}[section]
\newtheorem{lemma}[theorem]{Lemma}
\newtheorem{proposition}[theorem]{Proposition}
\newtheorem{corollary}[theorem]{Corollary}
\newtheorem*{theoremA}{Theorem A}
\newtheorem*{corollaryB}{Corollary B}
\newtheorem*{hypH}{Hypothesis $\mathrm{H}(\vartheta)$}
\newtheorem*{thmRS}{Theorem RS}
\theoremstyle{remark}
\newtheorem{remark}[theorem]{Remark}

\DeclareMathOperator{\pmin}{P^{-}}
\newcommand{\N}{N}
\newcommand{\J}{J_n}
\newcommand{\A}{\mathcal{A}}
\newcommand{\e}{\mathrm{e}}
\newcommand{\eof}[1]{e\!\left(#1\right)}
\newcommand{\thstar}{\vartheta^{**}}

\title{An Explicit Semiprime-Tail Bound in Square-Root Intervals\\
and a One-Hypothesis Reduction for $[4n^2-n,\,4n^2+n]$}
\author{Michael M. Ross\\
\textit{Independent Researcher}\\
\small \texttt{michaelmross@cantab.net}}

\begin{document}
\maketitle
\papermeta{Primary 11N35; Secondary 11N05, 11N25, 11N36, 11L07.}{%
Legendre's conjecture; primes in short intervals; linear sieve; Buchstab function; rough numbers; semiprimes; bilinear remainder terms; exponential sums with monomials; parity barrier; prime-detecting sieves.}

\begin{abstract}
Let $\J=[4n^2-n,\,4n^2+n]$ and $N=2n$. Every composite integer in $\J$ has a prime
factor at most $N$, so a Buchstab decomposition at $N^{\theta}$, $\theta>2/3$, writes
the number of primes in $\J$ as the count of $N^{\theta}$-rough integers minus a tail
$T(\theta)$ of semiprimes. We prove unconditionally that
$T(\theta)\le\bigl(C(\theta)+o(1)\bigr)u$ for every fixed $\theta\in(4/5,1)$, where
$u=n/\log N$ is the expected number of primes in $\J$ and
$C(\theta)=\frac83\log\frac{2-\theta}{5\theta-4}$, improving the linear-sieve bound
$(4+o(1))u$ for $\theta>0.8513$. The proof is a linear sieve over the cofactors of the
large prime factor, with the bilinear remainder controlled by the Robert--Sargos
exponential-sum estimate, and it holds for every interval $[x-y,\,x+y]$ with
$y\asymp x^{1/2}$. Consequently, for every fixed
$\vartheta>\vartheta^{**}=0.93930\ldots$, a lower bound at the Buchstab value for the
$N^{\vartheta}$-rough integers of $\J$ (Hypothesis $\mathrm{H}(\vartheta)$) implies
that $\J$ contains a prime for all sufficiently large $n$. We show that $\mathrm{H}(\vartheta)$ is a genuinely parity-breaking Type~II input and that the relevant bounded-weight Type~I remainders are already controlled unconditionally.
\end{abstract}

\section{Introduction and statement of results}\label{sec:intro}

Fix $n\ge 2$ and write throughout
\[
\N:=2n,\qquad x:=4n^2,\qquad \J:=[x-n,\,x+n],\qquad u:=\frac{n}{\log \N},
\]
and let $\A=\J\cap\mathbb{Z}$, so $|\A|=2n+1$. Since $x+n<(2n+1)^2$, every composite
$m\in\J$ has a prime factor $\le \N$ (Lemma~\ref{lem:structural}); hence, writing
$S(\A,z)=\#\{m\in\A:\pmin(m)>z\}$,
\begin{equation}\label{eq:pure-sieve}
\pi(\J):=\#\{p\in\J:\ p\ \text{prime}\}=S(\A,\N).
\end{equation}
Because $\J\subset\bigl((2n-1)^2,(2n+1)^2\bigr)$, a prime in $\J$ for every large $n$
would settle the quadratic-interval form of Legendre's conjecture. The broader
analytic setting is discussed in \cite{RossObstruction}.

A Buchstab decomposition of \eqref{eq:pure-sieve} at $\N^{\theta}$ with
$\theta\in(2/3,1)$ produces (Lemma~\ref{lem:buchstab})
\begin{equation}\label{eq:buchstab-intro}
\pi(\J)=S(\A,\N^{\theta})-T(\theta)+O(1),\qquad
T(\theta):=\#\{m\in\J:\ m=pq,\ \N^{\theta}<q\le \N<p\},
\end{equation}
with $p,q$ prime: for $\theta>2/3$ the discarded tail consists of semiprimes only.
The main result of this note is an unconditional upper bound for $T(\theta)$. The
classical bound is the Brun--Titchmarsh argument for the interval and uses no
exponential sums (the definition of $T(\theta)$ makes sense for every
$\theta\in(0,1)$).

\begin{proposition}[Classical tail bound]\label{prop:classical}
For every fixed $\theta\in[1/3,1)$, $T(\theta)\le(4+o(1))\,u$ as $n\to\infty$; and
$T(\theta)\ll_{\theta}u$ for every fixed $\theta\in(0,1)$.
\end{proposition}

The constant $4$ exceeds the Buchstab value $m(\vartheta)\le 1+\log 2$ of the main term
(Lemma~\ref{lem:constant}), so Proposition~\ref{prop:classical} cannot win the race in
Corollary~B below. The content of Theorem~A is the constant.

\begin{theoremA}
Fix $\theta\in(4/5,1)$ and put
\[
C(\theta):=\frac{8}{3}\,\log\frac{2-\theta}{5\theta-4}.
\]
Then, as $n\to\infty$,
\[
T(\theta)\le \bigl(C(\theta)+o(1)\bigr)\,u .
\]
\end{theoremA}

Since $C(\theta)<4$ exactly when $\theta>\theta_4:=(2+4e^{3/2})/(1+5e^{3/2})=0.85126\ldots$,
Theorem~A improves Proposition~\ref{prop:classical} on $(\theta_4,1)$.
Theorem~A is proved in Sections~\ref{sec:tail-setup}--\ref{sec:expsums} by a linear
sieve on the cofactors $q_p=m/p$, parametrized by the large prime $p$, at the per-slice
level
\begin{equation}\label{eq:profile}
D(a)=\N^{\zeta(a)},\qquad \zeta(a)=\frac{5a-4}{4}\qquad (q\asymp \N^{a},\ p\asymp \N^{2-a}),
\end{equation}
the bilinear sieve remainder being controlled through the Fouvry--Iwaniec reduction to
trilinear exponential sums \cite{FI} and the three-dimensional estimate of
Robert--Sargos \cite{RS}, after a Vaughan decomposition of the prime variable and an
Iwaniec-style bilinear splitting of the sieve modulus \cite{Iwaniec80}. The exponent
$(5a-4)/4$ is proved in Lemma~\ref{lem:coverage} and is optimal for these ingredients
(Proposition~\ref{prop:optimal}, Appendix~\ref{app:cert}). Nothing in the proof uses
the square center of $\J$, and Theorem~A holds for every interval $[x-y,\,x+y]$ with
$y\asymp x^{1/2}$ (Section~\ref{sec:general}).

The main term $S(\A,\N^{\vartheta})$ is governed by the following hypothesis, in which
$\omega$ denotes the Buchstab function and we use the classical identity
$s\,\omega(s)=1+\log(s-1)$ for $s\in[2,3]$.

\begin{hypH}
Fix $\vartheta\in(2/3,1)$. As $n\to\infty$,
\[
S(\A,\N^{\vartheta})\ \ge\ \bigl(1-o(1)\bigr)\,m(\vartheta)\,u,
\qquad m(\vartheta):=1+\log\frac{2-\vartheta}{\vartheta}.
\]
\end{hypH}

The quantity $m(\vartheta)u=\frac{2}{\vartheta}\,\omega(2/\vartheta)\,u$ is exactly
the Buchstab prediction for the number of $\N^{\vartheta}$-rough integers in an
interval of length $2n$ near $x=\N^2$ (Lemma~\ref{lem:constant}); only the lower bound
is assumed, and it is assumed with no loss. Hypothesis $\mathrm{H}(\vartheta)$ is open
for every $\vartheta$ in the stated range. Section~\ref{sec:hyp} characterizes it:
the Type~I remainders of the relevant sieve, against arbitrary bounded weights, are
controlled unconditionally at level $\N^{1-\varepsilon}$ (Lemma~\ref{lem:remainder}),
while the linear sieve, at any level of distribution below $x$, falls short of the
required bound by a positive proportion of $u$ (Proposition~\ref{prop:parity}).
$\mathrm{H}(\vartheta)$ is therefore a parity-breaking hypothesis, and the input it
requires is Type~II information for the sequence $\mathbf{1}_{\J}$
(Remark~\ref{rem:requires}).

\begin{corollaryB}[Conditional prime consequence]
Let $\thstar=0.939305\ldots$ be the unique root in $(4/5,1)$ of
\begin{equation}\label{eq:threshold}
1+\log\frac{2-\vartheta}{\vartheta}\ =\ \frac83\,\log\frac{2-\vartheta}{5\vartheta-4}.
\end{equation}
Assume Hypothesis $\mathrm{H}(\vartheta)$ for some fixed $\vartheta\in(\thstar,1)$.
Then, as $n\to\infty$,
\[
\pi(\J)\ \ge\ \bigl(c(\vartheta)-o(1)\bigr)\,u,\qquad
c(\vartheta):=m(\vartheta)-C(\vartheta)>0 .
\]
In particular, for all sufficiently large $n$ the interval $[4n^2-n,\,4n^2+n]$
contains a prime.
\end{corollaryB}

Corollary~B follows immediately from Theorem~A and \eqref{eq:buchstab-intro}; the threshold is checked in Section~\ref{sec:race}.

\subsection*{Notation}
$p,q$ denote primes; $\pmin(m)$ the least prime factor of $m$ ($\pmin(1)=\infty$);
$P(z)=\prod_{p<z}p$; $\psi(t)=t-\lfloor t\rfloor-\tfrac12$;
$e(t)=e^{2\pi i t}$; $L:=\log \N$. We write $m\sim M$ for $M<m\le 2M$, and
$f\ll g$ with implied constants absolute unless subscripted. Throughout,
$\varepsilon$ denotes a sufficiently small positive constant, not the same at each
occurrence, and $o(1)$ refers to $n\to\infty$ with all other parameters fixed.
$F$ and $f$ denote the upper and lower functions of the linear sieve,
$F(s)=2e^{\gamma_0}/s$ for $1\le s\le 3$, $f(s)=0$ for $s\le 2$ and
$f(s)=2e^{\gamma_0}\log(s-1)/s$ for $2\le s\le 4$, $\gamma_0$ Euler's constant
\cite{FriedlanderIwaniec}. For a set $S$ we write $\mathbf{1}_S$ for its indicator function. We identify $\A$ with the sequence $a_m=\mathbf{1}_{\J}(m)$, $m\ge 1$, in the sense of \cite{FriedlanderIwaniec}, and use whichever form reads more clearly.

\section{Bookkeeping}\label{sec:bookkeeping}

\begin{lemma}\label{lem:structural}
Every composite $m\in\J$ has a prime factor $\le \N$, and \eqref{eq:pure-sieve} holds.
\end{lemma}

\begin{proof}
If $m\in\J$ is composite then $\pmin(m)\le\sqrt m\le\sqrt{4n^2+n}<2n+1$, so
$\pmin(m)\le 2n=\N$. Every $m\in\J$ exceeds $\N$, so the $\N$-rough elements of $\A$
are exactly the primes in $\J$.
\end{proof}

\begin{lemma}[Buchstab step]\label{lem:buchstab}
For fixed $\theta\in(2/3,1)$, identity \eqref{eq:buchstab-intro} holds with
$T(\theta)$ counting semiprimes $m=pq$, $\N^{\theta}<q\le\N<p$.
\end{lemma}

\begin{proof}
Buchstab's identity gives
$S(\A,\N)=S(\A,\N^{\theta})-\sum_{\N^{\theta}<q\le \N}S(\A_q,q)$,
where $\A_q=\{m\in\A: q\mid m\}$ and $S(\A_q,q)$ counts $m=q\,m'$ with
$\pmin(m')\ge q$. Here $m'\le (x+n)/\N^{\theta}<\N^{2-\theta}(1+o(1))$, and
$\pmin(m')>\N^{\theta}$ with $2\theta>2-\theta$ forces $m'=1$ or $m'$ prime.
The case $m'=1$ is impossible ($q\le\N<\min\J$). If $m'=p$ then
$p=m/q\ge (x-n)/\N>\N-1$, and $p=\N$ is impossible ($\N$ even), so $p>\N\ge q$.
The case $p=q$ would force $m=q^2\in\J$, whence $q\in(2n-1,2n+1)$, i.e.\ $q=2n$,
not prime. Finally $S(\A,\N)=\pi(\J)$ by Lemma~\ref{lem:structural}; the $O(1)$
absorbs boundary conventions at $q=\N^{\theta}$.
\end{proof}

\begin{proof}[Proof of Proposition~\ref{prop:classical}]
Every $m$ counted by $T(\theta)$ has both prime factors above $\N^{\theta}$, so
$T(\theta)\le S(\A,\N^{\theta})$. For the interval $\A$ one has
$|\A_d|=(2n+1)/d+r_d$ with $|r_d|\le 1$, so the sieve axioms hold with $X=2n+1$,
$g(d)=1/d$ and $\sum_{d\le D}|r_d|\le D$. The upper-bound linear sieve at level
$D=\N^{1-\varepsilon}$ with $z=\N^{\theta}$ gives
$S(\A,z)\le (F(s)+o(1))(2n+1)V(z)+\N^{1-\varepsilon}$, $s=(1-\varepsilon)/\theta$.
For $\theta\in[1/3,1)$ and $\varepsilon$ small, $s\in[1,3]$, where $F(s)=2e^{\gamma_0}/s$
and hence
\[
(2n+1)V(z)F(s)=(2n+1)(1+o(1))\frac{2}{s\log z}=(2n+1)(1+o(1))\frac{2}{\log D}
=\bigl(4+O(\varepsilon)+o(1)\bigr)u .
\]
Let $\varepsilon\to 0$ along a sequence. For $\theta\in(0,1/3)$ the same bound with
$s>3$ gives $S(\A,z)\ll u/\theta$, since $F$ is bounded.
\end{proof}

\begin{lemma}[Main-term constant]\label{lem:constant}
The Buchstab prediction for $S(\A,\N^{\vartheta})$, $\vartheta\in[2/3,1)$, is
$m(\vartheta)\,u$ with $m(\vartheta)=1+\log\frac{2-\vartheta}{\vartheta}$.
\end{lemma}

\begin{proof}
The density of $z$-rough integers near $x$ is $\omega(s)/\log z$ with
$s=\log x/\log z$ \cite[III.6]{Tenenbaum}. With $z=\N^{\vartheta}$ and $x=\N^{2}(1+o(1))$
we get $s=2/\vartheta\in[2,3]$ and prediction
$2n\cdot \omega(2/\vartheta)/(\vartheta L)=\tfrac{2}{\vartheta}\omega(2/\vartheta)\,u$.
On $[2,3]$, $(s\omega(s))'=\omega(s-1)=1/(s-1)$ and $2\omega(2)=1$, so
$s\omega(s)=1+\log(s-1)$; substituting $s=2/\vartheta$ gives the claim.
\end{proof}

\section{The tail: swap and sieve}\label{sec:tail-setup}

\begin{lemma}\label{lem:swap}
For each prime $p>\N$ there is at most one multiple of $p$ in $\J$; write it $p\,q_p$
when it exists. Then
\[
T(\theta)\ \le\ \#\{p\ \mathrm{prime},\ \N<p\le \N^{2-\theta}(1+o(1)):\ q_p\ \text{exists and is prime}\}.
\]
\end{lemma}

\begin{proof}
The interval $[(x-n)/p,(x+n)/p]$ has length $2n/p<1$. If $m=pq\in\J$ with
$q>\N^{\theta}$ then $p=m/q<(x+n)/\N^{\theta}$.
\end{proof}

We bound the right side by an upper-bound linear sieve on the family of cofactors
$q_p$, slice by slice. Fix a dyadic block $P\subset(\N,\N^{2-\theta}]$, write
$P=\N^{\gamma}$, $\gamma\in[1,2-\theta]$, and set $a:=2-\gamma$, so
$q_p\asymp \N^{a}$. Let $z:=\N^{\zeta/2}$ and relax primality of $q_p$ to
$\pmin(q_p)>z$; this is a relaxation since $q_p\asymp \N^{a}\ge \N^{4/5}>z$.
For squarefree $e\mid P(z)$ put
\[
\A_e(P):=\#\{(p,j):\ p\sim P\ \mathrm{prime},\ e p j\in\J\}
= X_P\,\frac1e + r_e(P),\qquad X_P:=(2n+1)\sum_{p\sim P}\frac1p ,
\]
where, with $\psi(t)=t-\lfloor t\rfloor-\tfrac12$,
\begin{equation}\label{eq:remainder-def}
r_e(P)=\sum_{p\sim P}\Bigl(\psi\Bigl(\frac{x-n-1}{ep}\Bigr)-\psi\Bigl(\frac{x+n}{ep}\Bigr)\Bigr).
\end{equation}
The sharp block $p\sim P$ may be replaced throughout by a smooth one, $p$ weighted by
$g(p/P)$ with $g\in C_c^{\infty}([1/2,2])$ nonnegative, with $X_P:=(2n+1)\sum_p g(p/P)/p$;
the sieve below applies to nonnegative weights verbatim, and
Section~\ref{sec:expsums} uses this freedom (Lemma~\ref{lem:separation}).

\begin{proposition}[Slice bound]\label{prop:slice}
Let $0<\zeta\le 1/2$ and $D:=\N^{\zeta}$, and suppose that for every choice of
coefficients $(\xi_e)$ with $|\xi_e|\le 1$ supported on squarefree $e\le D$ and
representable in Iwaniec's bilinear form at any prescribed factorization $D=D_1D_2$,
\begin{equation}\label{eq:interface}
\Bigl|\sum_{e}\xi_e\, r_e(P)\Bigr|\ \ll\ \N^{1-\varepsilon}.
\end{equation}
Then
\[
\#\{p\sim P:\ q_p\ \text{exists},\ \pmin(q_p)>z\}
\ \le\ \bigl(1+o(1)\bigr)\,\frac{2X_P}{\zeta\, L}.
\]
\end{proposition}

\begin{proof}
Apply the upper linear sieve in Iwaniec's form \cite{Iwaniec80},
\cite[Ch.~12]{FriedlanderIwaniec}: with level $D$ and $z=D^{1/2}$ (so $s=2$),
\[
\sum_{p\sim P}\mathbf{1}\{\pmin(q_p)>z\}
\le X_P\,V(z)\bigl(F(2)+o(1)\bigr)+R,
\qquad V(z)F(2)=\frac{2e^{-\gamma_0}\cdot e^{\gamma_0}}{\log D}\bigl(1+o(1)\bigr),
\]
and the error $R$ is a bounded combination of
$O_\varepsilon(1)$ bilinear forms
\[
\sum_{e_1\le D_1}\sum_{e_2\le D_2} a_{e_1}b_{e_2}\,r_{e_1e_2}(P),
\qquad |a|,|b|\le 1,
\]
at any prescribed factorization $D_1D_2=D$; each is
$\ll \N^{1-\varepsilon}=o(X_P/L)$ by \eqref{eq:interface}. The weights $\lambda^{+}$
are fixed once for all; the factorization enters only through the representation of
the remainder, which Iwaniec's theorem supplies for every factorization $D_1D_2=D$
simultaneously. Since each of the blocks into which Section~\ref{sec:expsums} splits
the remainder is a linear functional of $\lambda^{+}$, the factorization may be chosen
block by block.
\end{proof}

Summing Proposition~\ref{prop:slice} over dyadic blocks, with the slice-dependent
level $\zeta=\zeta(a)$ of \eqref{eq:profile} and Mertens'
$\sum_{p\sim P}1/p=(1+o(1))\log 2/\log P$, yields the Riemann sum of
\begin{equation}\label{eq:tail-integral}
T(\theta)\ \le\ (1+o(1))\,u\int_{\theta}^{1}\frac{4\,da}{(2-a)\,\zeta(a)}\ +\ O(\N^{1-\varepsilon}L),
\end{equation}
so Theorem~A reduces to establishing \eqref{eq:interface} at level
$\zeta(a)=(5a-4)/4-\varepsilon_0$ for each fixed $\varepsilon_0>0$, plus the closed-form
evaluation of the integral (Section~\ref{sec:expsums-assemble}). Everything now rests
on the exponential-sum estimate \eqref{eq:interface}, and Section~\ref{sec:expsums}
is devoted to it.

\section{The tail: exponential sums}\label{sec:expsums}

Throughout this section fix a slice: $P=\N^{\gamma}$ with $\gamma\in[1,6/5)$,
a sieve modulus $e=e_1e_2$ with $e_1\sim E_1$, $e_2\sim E_2$ in independent dyadic
ranges, $E:=E_1E_2=\N^{\beta'}$ with $\beta'\le\beta:=\zeta-\varepsilon_0$, and write
\[
\Sigma:=\beta+\gamma\ \Bigl(\le \tfrac{6-\gamma}{4}\le\tfrac54\Bigr),
\qquad \xi:=\eta+2-\Sigma ,
\]
where $\N^{\eta}$ will denote a dyadic frequency block. All exponents are with respect
to base $\N$, and $x\asymp \N^2$.

\subsection{From remainders to trilinear sums}\label{sec:pipeline}

\begin{lemma}[Pipeline; after {\cite[Lemma~9]{FI}}]\label{lem:pipeline}
Let $E_1,E_2,P\ge 1$ with $E_1E_2P\le x$, write $E:=E_1E_2$, let
$|c_{e_1}|,|c'_{e_2}|,|b_p|\le 1$, and set $y:=2n+1$, $H:=4EP\,y^{-1}x^{3\varepsilon}$.
If for every $z'\in[x-2y,\,x+2y]$ and every $1\le H'\le H$,
\begin{equation}\label{eq:trilinear-req}
\Bigl|\sum_{h\le H'}\sum_{e_1\sim E_1}\sum_{e_2\sim E_2}\sum_{p\sim P}
c_{e_1}c'_{e_2}\,b_p\; \eof{\frac{h z'}{e_1e_2 p}}\Bigr|
\ \le\ EP\,x^{-2\varepsilon},
\end{equation}
then
\[
\Bigl|\sum_{e_1\sim E_1}\sum_{e_2\sim E_2}\sum_{p\sim P}c_{e_1}c'_{e_2}\,b_p\,r_{e_1e_2p}\Bigr|
\ \ll\ y\,x^{-\varepsilon},
\qquad
r_d:=\Bigl\lfloor\frac{x+n}{d}\Bigr\rfloor-\Bigl\lfloor\frac{x-n-1}{d}\Bigr\rfloor-\frac{y}{d}.
\]
\end{lemma}

\begin{proof}
This is \cite[Lemma~9]{FI} with $(x,y)$ there replaced by $(x+n,\,2n+1)$ here. The
argument is insensitive to the ranges of the outer variables, which enter only through
the truncation height $H\asymp \max(e p)\,y^{-1}x^{3\varepsilon}$, and here $e\le 4E$.
In brief: sandwich the interval indicator between smooth functions $f_{\pm}$ of total
variation $O(1)$, supported within $y\,x^{-2\varepsilon}$ of $\J$, with
$f_\pm^{(j)}\ll (y x^{-2\varepsilon})^{-j}$; Poisson summation in the inner variable
converts $\sum_j f_\pm(epj)$ into $\widehat{f_\pm}(0)/(ep)$ plus frequencies $h\ne 0$.
Frequencies $|h|>H$ are negligible after $j$ integrations by parts with
$j\gg 1/\varepsilon$, since $|h|/(ep)\cdot y x^{-2\varepsilon}\ge x^{\varepsilon}$
there. The main terms reassemble $y/(ep)$ up to $O(y x^{-2\varepsilon})$ total, and the
surviving frequencies are bounded by suprema of the sums \eqref{eq:trilinear-req}
after removing the smooth weights $\widehat{f_\pm}(h/ep)/(ep)$ by partial summation
in each variable (the weights are monotone in dyadic ranges), at the cost of the
factor $x^{\varepsilon}$ absorbed in \eqref{eq:trilinear-req}.
\end{proof}

By dyadic decompositions (a factor $O(L^4)$, absorbed in $x^{-\varepsilon}$),
\eqref{eq:trilinear-req} follows from: \emph{for every dyadic
$H_0=\N^{\eta}$, $0\le \eta\le \Sigma-1+O(\varepsilon)$, the block sum with $h\sim H_0$
is at most $EP\,x^{-3\varepsilon}$}. Blocks with $\beta'+\gamma\le 1-O(\varepsilon)$
satisfy $H<1$ and require nothing. We henceforth fix a block and prove the required
bound.

\subsection{The obstruction to a direct estimate}\label{sec:obstruction}

Estimates for sums \eqref{eq:trilinear-req} with two free coefficient slots ---
\cite[Thm.~6, Thm.~7]{FI}, or Theorem~RS below applied with the pair $(e,p)$ ---
cannot be applied with the prime $p$ occupying a slot as-is. The packaged sieve form
\cite[Thm.~7]{FI} requires \emph{both} coefficient ranges below the interval length
(their conditions (7.4)--(7.6) force $M<y\,x^{-\varepsilon'}$ and $N^6<My^{7}x^{-3}$),
while $p\sim \N^{\gamma}\ge \N=y(1+o(1))$. The same failure is visible directly in
Theorem~RS: with $p$ undecomposed the admissible standalone exponents are
$m_1\in[0,\beta']\cup[\gamma,\gamma+\beta']$, whereas the window of
Lemma~\ref{lem:window} below is $[2\eta,\Sigma-3\eta]$, which misses both components
as soon as $\eta>\beta'/2$ --- and $\eta$ ranges up to $\Sigma-1\ge\gamma-1$, with
$\gamma>1$ on all interior slices. A decomposition of the prime variable is therefore
forced, and the entire question becomes which factorizations the resulting pieces
admit. This coefficient-restructuring problem is analyzed in \cite{RossObstruction}; the
remainder of the section settles it for the present cofactor sums.

\subsection{Decomposition and grouping}\label{sec:decomp}

Apply Vaughan's identity \cite{Vaughan} to the prime variable with cutoffs
$U=V=P^{1/3}$: the indicator of primes $p\sim P$ (weighted by $\log p$, removed by
partial summation) decomposes into $O(1)$ sums of two shapes,
\[
\textbf{(I)}\ \sum_{m_v\le P^{2/3}}\alpha_{m_v}\sum_{k}(\ast)_k,\qquad
\textbf{(II)}\ \sum_{P^{1/3}<m_v\le P^{1/2}}\ \sum_{k}\alpha_{m_v}\beta_k ,
\]
with $m_vk\sim P$, where in type~I the inner weight $(\ast)_k\in\{1,\log k\}$ is
smooth and in type~II both coefficient sequences are arbitrary with divisor-bounded
size; all coefficients are $\ll \N^{o(1)}$, absorbed into $x^{\varepsilon}$-slack.
Writing $m_v=\N^{\mu}$, the pieces carry
\[
\textbf{(I)}:\ \mu\in[0,\tfrac{2\gamma}{3}],\qquad
\textbf{(II)}:\ \mu\in[\tfrac{\gamma}{3},\tfrac{\gamma}{2}]
\]
(after choosing the shorter type-II side). Simultaneously, by
Proposition~\ref{prop:slice} the sieve coefficients may be prescribed in bilinear
form $e=e_1e_2$ with $e_1\le \N^{\beta-s}$, $e_2\le \N^{s}$ for \emph{any}
$s\in[0,\beta]$, chosen per block.

The estimates below (Theorem~RS and Lemma~\ref{lem:elementary}) require their groups
of variables to range over independent dyadic boxes with coefficients in product
form. The following lemma supplies that form.

\begin{lemma}[Separation of variables]\label{lem:separation}
Each block of \eqref{eq:trilinear-req} may be reduced, at a cost
$\N^{O(\varepsilon)}$ absorbed in the available slack and $O(L)$ additional
blocks, to a sum in which:
the sieve variables $e_1\sim E_1$ and $e_2\sim E_2$ range independently;
the Vaughan variables have product-form coefficients; and every group used as
a joint or standalone variable in Theorem~RS ranges over a dyadic interval
with arbitrary bounded coefficients. In a type-I piece in which $k$ remains
free, its weight is smooth with $j$-th derivative $\ll_j K^{-j}$.
\end{lemma}

\begin{proof}
There are two apparent obstacles to product structure. First, the sieve
condition may couple $e_1$ and $e_2$ through their product $e=e_1e_2$.
Second, after Vaughan's identity the condition $p\sim P$ becomes
$m_vk\sim P$.

The first obstacle is avoided by construction: throughout we keep
$e_1\sim E_1$ and $e_2\sim E_2$ in independent dyadic ranges. This is the
reason Lemma~\ref{lem:pipeline} allows $e=e_1e_2$ to range over the
factor-$4$ interval $(E,4E]$.

For the Vaughan variables, use the smooth $p$-partition introduced in
Section~\ref{sec:tail-setup}. Thus the sharp condition $m_vk\sim P$ is
replaced by the smooth weight
\[
G(m_vk):=\frac{g(m_vk/P)}{\log(m_vk)},
\qquad g\in C_c^\infty([1/2,2]),
\]
at a cost of $O(L)$ blocks.

In a type-II piece, and in a type-I piece when $k$ is to occupy a
standalone Robert--Sargos slot, separate this weight by Mellin inversion:
\[
G(mk)=\frac{1}{2\pi i}\int_{(0)}
   \widehat G(s)m^{-s}k^{-s}\,ds .
\]
Since $G$ is smooth and supported on $[P/2,2P]$,
\[
\widehat G(i\tau)\ll_j (1+|\tau|)^{-j}.
\]
The integral may therefore be truncated at $|\tau|\le \N^\varepsilon$
with negligible error. For fixed $\tau$, the factors $m^{-i\tau}$ and
$k^{-i\tau}$ have modulus $1$ and may be absorbed into the coefficients.
The remaining $\tau$-integration costs at most $\N^\varepsilon$, which is
absorbed by the slack in \eqref{eq:trilinear-req}.

In a type-I piece with $k$ left free, no Mellin separation is needed:
we retain $G(m_vk)$ as a smooth weight in $k$. Its $j$-th derivative is
$\ll_j K^{-j}$, and Lemma~\ref{lem:elementary} remains valid after partial
summation.

Finally, a Robert--Sargos slot may consist of a product of several variables,
such as $e_1e_2m_v$ or $m_ve_2$. Collapse such a product to a single variable,
say $\nu$. Its coefficient is the sum over the factorizations of $\nu$ and
is therefore divisor-bounded; this is absorbed into the
$\N^{O(\varepsilon)}$ loss. Splitting the resulting range into dyadic
intervals completes the reduction.
\end{proof}

Each block of \eqref{eq:trilinear-req} is now a sum over
$(h;e_1,e_2,m_v,k)$ with phase $\eof{hz'/(e_1e_2m_vk)}$. To estimate it we designate
one group of variables as the \emph{standalone} slot $w$ (of exponent $m_1$) and the
rest join $h$ in a \emph{joint} slot; the available standalone exponents are
\begin{equation}\label{eq:available}
m_1\in \mathcal{A}(\mu):=[0,\beta]\ \cup\ [\mu,\,\mu+\beta]\ \cup\ [\gamma-\mu,\,\gamma-\mu+\beta],
\end{equation}
realized by $w\in\{e_2,\ m_ve_2,\ k e_2\}$ with $s\in[0,\beta]$ tuned continuously.

\subsection{The engine}\label{sec:engine}

\begin{thmRS}[{Robert--Sargos \cite[Thm.~1]{RS}}]
Let $\alpha,\beta_0,\gamma_0$ be fixed nonzero reals with
$\alpha(\alpha-1)\beta_0\gamma_0\ne 0$, let $H,N_1,M_1\ge 1$, $X\ge 1$, and let
$a(h,\nu)$, $b(w)$ be arbitrary complex coefficients of modulus $\le 1$. Then
\begin{multline*}
\sum_{h\sim H}\sum_{\nu\sim N_1}a(h,\nu)\sum_{w\sim M_1}b(w)\,
\eof{X\,\frac{h^{\beta_0}\nu^{\gamma_0}w^{\alpha}}{H^{\beta_0}N_1^{\gamma_0}M_1^{\alpha}}}\\
\ll_{\varepsilon}\ (HN_1M_1)^{1+\varepsilon}
\Bigl(\Bigl(\frac{X}{HN_1M_1^{2}}\Bigr)^{1/4}+\frac{1}{(HN_1)^{1/4}}
+\frac{1}{M_1^{1/2}}+\frac{1}{X^{1/2}}\Bigr).
\end{multline*}
\end{thmRS}

We apply this with $(\beta_0,\gamma_0,\alpha)=(1,-1,-1)$ (so
$\alpha(\alpha-1)\beta_0\gamma_0=-2\ne0$): $h$ and the joint group occupy $(h,\nu)$,
the standalone group occupies $w$, and $X=\N^{\xi}$ is the phase size. Incomplete and
maximal ranges arising from partial summation are standard
\cite[Lemma~2]{RS}.

\begin{lemma}[Admissibility window]\label{lem:window}
Let the block have $h\sim\N^{\eta}$, standalone exponent $m_1$, joint exponent
$\Sigma-m_1$, with $0\le\eta\le\Sigma-1$ and $\Sigma\le 5/4$. The Theorem~RS bound is
$\le \N^{\Sigma-c\varepsilon}$ if and only if
\[
2\eta\ \le\ m_1\ \le\ \Sigma-3\eta
\]
(up to $O(\varepsilon)$ in each inequality).
\end{lemma}

\begin{proof}
The trivial size is $\N^{\eta+\Sigma}$, so the requirement is that each of the four
bracket terms be $\le \N^{-\eta-c\varepsilon}$. In exponents, with $n_1=\Sigma-m_1$:
\[
\text{(i)}\ \frac{\xi-\eta-n_1-2m_1}{4}=\frac{2-2\Sigma-m_1}{4}\le-\eta
\iff m_1\ge 4\eta+2-2\Sigma;
\]
\[
\text{(ii)}\ -\frac{\eta+n_1}{4}\le-\eta\iff m_1\le \Sigma-3\eta;\qquad
\text{(iii)}\ -\frac{m_1}{2}\le-\eta\iff m_1\ge 2\eta;
\]
\[
\text{(iv)}\ -\frac{\xi}{2}\le-\eta\iff \eta\le 2-\Sigma .
\]
Condition (iv) holds since $\eta\le\Sigma-1\le\frac14\le 2-\Sigma$ for
$\Sigma\le\frac54$; condition (i) is implied by (iii) since
$4\eta+2-2\Sigma\le 2\eta\iff\eta\le\Sigma-1$. This leaves (ii) and (iii).
\end{proof}

\subsection{The elementary bound for type-I pieces}\label{sec:elementary}

\begin{lemma}\label{lem:elementary}
Consider a type-I block: modulus $c=e_1e_2m_v\sim \N^{\chi}$, $\chi=\beta'+\mu$,
carrying arbitrary bounded coefficients, and smooth inner variable
$k\sim K=\N^{\kappa}$, $\kappa=\gamma-\mu$. If
\[
3\eta+2\mu\ \le\ \beta'+3\gamma-2 - c\varepsilon,
\]
then the block sum is $\le \N^{\Sigma'-c'\varepsilon}$, $\Sigma'=\beta'+\gamma$.
\end{lemma}

\begin{proof}
Sum absolutely over $h$ and $c$ and estimate the $k$-sum. The phase
$f(k)=hz'/(ck)$ has $|f'|\asymp \Lambda:=\N^{\lambda}$,
$\lambda=\eta+2-\chi-2\kappa$, monotone with $|f''|\asymp\Lambda/K$. By the
exponent pair $(\tfrac12,\tfrac12)$ (van der Corput's second-derivative test,
\cite[Thm.~2.2]{GK}), together with the Kuzmin--Landau complement,
\[
\Bigl|\sum_{k\sim K}\eof{f(k)}\Bigr|\ \ll\ \Lambda^{1/2}K^{1/2}+\Lambda^{-1}.
\]
The first term contributes exponent
$\eta+\chi+\tfrac{\lambda}{2}+\tfrac{\kappa}{2}
=\tfrac{3\eta}{2}+\tfrac{\beta'}{2}+\mu+1-\tfrac{\gamma}{2}$, and the requirement
$\le \Sigma'$ is the displayed inequality. The second contributes
$\eta+\chi+(\chi+2\kappa-\eta-2)=2\Sigma'-2\le \Sigma'-\tfrac34$ since
$\Sigma'\le\frac54$.
\end{proof}

\subsection{The coverage lemma}\label{sec:coverage}

\begin{lemma}[Coverage]\label{lem:coverage}
Let $\gamma\in[1,6/5)$ and
\[
0<\beta\ \le\ \frac{6-5\gamma}{4}-\varepsilon_0,\qquad \Sigma=\beta+\gamma .
\]
Then for every Vaughan piece and every $\eta\in[0,\Sigma-1]$, at least one admissible
configuration exists:
\begin{itemize}
\item[\emph{(II)}] for type-II pieces $\mu\in[\gamma/3,\gamma/2]$: Theorem~RS with
standalone exponent $m_1=\max(\mu,2\eta)\in[\mu,\mu+\beta]$;
\item[\emph{(I)}] for type-I pieces $\mu\in[0,2\gamma/3]$: Lemma~\ref{lem:elementary}
if $3\eta+2\mu\le\beta+3\gamma-2$; otherwise Theorem~RS with
$m_1=\max(\mu,2\eta)\in[\mu,\mu+\beta]$ if $\mu\le\Sigma-3\eta$; otherwise Theorem~RS
with $m_1=\max(\gamma-\mu,2\eta)\in[\gamma-\mu,\gamma-\mu+\beta]$.
\end{itemize}
Consequently every block of \eqref{eq:trilinear-req} is $\le EP\,x^{-3\varepsilon}$,
and \eqref{eq:interface} holds at level $\zeta=\beta$.
\end{lemma}

\begin{proof}
All inequalities below hold with linear slack in $\varepsilon_0$; we suppress the
$\varepsilon$'s. Note first $\Sigma\le\gamma+\frac{6-5\gamma}{4}=\frac{6-\gamma}{4}
\le\frac54$, so Lemma~\ref{lem:window} applies, and $\eta\le\Sigma-1\le\Sigma/5$,
so $2\eta\le\Sigma-3\eta$: the window is nonempty. Five verifications remain.

\smallskip
\emph{Type II.} (V1) Upper edge: for all $\eta\le\Sigma-1$,
\[
\mu\ \le\ \frac{\gamma}{2}\ \le\ 3-2\Sigma\ \le\ \Sigma-3\eta ,
\]
where $\gamma/2\le 3-2\Sigma$ is exactly $4\beta+5\gamma\le 6$. (V2) Lower edge:
$m_1\le\mu+\beta$ requires $2\eta\le\mu+\beta$, and for all $\eta\le\Sigma-1$,
\[
2(\Sigma-1)\le \frac{\gamma}{3}+\beta
\iff \beta\le 2-\frac{5\gamma}{3},
\]
which follows from $\beta\le\frac{6-5\gamma}{4}$ on $\gamma\le\frac65$ since
$\frac{6-5\gamma}{4}\le 2-\frac{5\gamma}{3}\iff 5\gamma\le 6$. Thus
$m_1=\max(\mu,2\eta)$ lies in $[\mu,\mu+\beta]\cap[2\eta,\Sigma-3\eta]$.

\smallskip
\emph{Type I.} At a given $\eta$, the three tools cover the $\mu$-intervals
\[
\Bigl[0,\ \mu_E(\eta)\Bigr]\ \cup\
\Bigl[\max(0,2\eta-\beta),\ \Sigma-3\eta\Bigr]\ \cup\
\Bigl[3\eta-\beta,\ \gamma+\beta-2\eta\Bigr],
\qquad \mu_E(\eta):=\frac{\beta+3\gamma-2-3\eta}{2},
\]
the second and third by the same window computation as in (V1)--(V2) applied to
$m_1=\max(\mu,2\eta)$ and $m_1=\max(\gamma-\mu,2\eta)$ respectively. Their union
contains $[0,2\gamma/3]$ provided the following three inequalities hold for all
$\eta\le \Sigma-1$; each is checked at the worst point $\eta=\Sigma-1$:
\[
\text{(V3)}\quad 2\eta-\beta\ \le\ \mu_E(\eta)
\iff 7\eta\le 3\Sigma-2
\ \Longleftarrow\ 7(\Sigma-1)\le 3\Sigma-2\iff \Sigma\le\tfrac54 ;
\]
\[
\text{(V4)}\quad 3\eta-\beta\ \le\ \Sigma-3\eta
\iff 6\eta\le\Sigma+\beta
\ \Longleftarrow\ 6(\Sigma-1)\le\Sigma+\beta\iff 4\beta+5\gamma\le 6 ;
\]
\[
\text{(V5)}\quad \frac{2\gamma}{3}\ \le\ \gamma+\beta-2\eta
\iff 2\eta\le\beta+\frac{\gamma}{3}
\ \Longleftarrow\ \beta\le 2-\frac{5\gamma}{3}\quad\text{as in (V2)} .
\]
This completes the case analysis: (V1) and (V4) are the facet $4\beta+5\gamma\le6$,
(V2) and (V5) are implied by it on $\gamma\in[1,6/5]$, and (V3) by
$\Sigma\le\frac54$, also implied. Each inequality is tightest at $\eta=\Sigma-1$ and
monotone in $\beta'\le\beta$ in the right direction, so the verifications survive
every sub-block $\beta'\le\beta$.
\end{proof}

\subsection{Optimality within the method}\label{sec:optimal}

\begin{proposition}[Optimality within the method]\label{prop:optimal}
Let $\gamma\in[1,6/5)$ and $\beta>\frac{6-5\gamma}{4}$. Then the balanced type-II
block $\mu=\gamma/2$ at the top frequency $\eta=\Sigma-1$ admits no admissible
configuration from the tools of Lemma~\ref{lem:coverage}. Consequently the level
$\zeta(\gamma)=\frac{6-5\gamma}{4}$ is exact for that toolkit: Theorem~RS, the
van der Corput and Kuzmin--Landau bounds, Vaughan cutoffs at $P^{1/3}$, and
arbitrary bilinear splitting of the sieve modulus.
\end{proposition}

\begin{proof}
The elementary bound needs a smooth inner variable, which a type-II piece lacks. For
Theorem~RS the window of Lemma~\ref{lem:window} at $\eta=\Sigma-1$ is
$[2\Sigma-2,\,3-2\Sigma]$, and $\beta>\frac{6-5\gamma}{4}$ gives $3-2\Sigma<\gamma/2$,
so the families $[\mu,\mu+\beta]$ and $[\gamma-\mu,\gamma-\mu+\beta]$ of
\eqref{eq:available}, both of which start at $\gamma/2$, lie above the window, while
the family $[0,\beta]$ would need $\beta\ge 2\Sigma-2$, i.e.\ $\beta\le 2-2\gamma\le 0$.
Lemma~\ref{lem:coverage} attains the level for every
$\beta\le\frac{6-5\gamma}{4}-\varepsilon_0$.
\end{proof}

The binding constraint is the balanced type-II block at top frequency. Appendix~\ref{app:cert} verifies in exact rational arithmetic that adding the further Fouvry--Iwaniec groupings and the exponent pair $(\tfrac16,\tfrac23)$ does not improve the profile.

\subsection{Proof of Theorem A}\label{sec:expsums-assemble}

Fix $\theta\in(4/5,1)$ and $\varepsilon_0>0$. For each dyadic slice
$P=\N^{\gamma}$, $\gamma\in[1,2-\theta]\subset[1,6/5)$, take the level
$\zeta=\zeta(a)-\varepsilon_0=\frac{5a-4}{4}-\varepsilon_0$ with $a=2-\gamma$.
Lemma~\ref{lem:coverage} verifies the hypothesis of Lemma~\ref{lem:pipeline} for
every sub-block $\beta'\le\beta$, hence \eqref{eq:interface};
Proposition~\ref{prop:slice} then gives the slice bound,
and \eqref{eq:tail-integral} gives, using
$5(2-a)+(5a-4)=6$ for the partial fractions,
\[
\int_{\theta}^{1}\frac{16\,da}{(2-a)(5a-4)}
=\frac{16}{6}\Bigl[\log\frac{5a-4}{2-a}\Bigr]_{\theta}^{1}
=\frac83\,\log\frac{2-\theta}{5\theta-4}\ =\ C(\theta).
\]
Letting $\varepsilon_0\to0$ along a sequence (the constant depends continuously on
$\varepsilon_0$) completes the proof. \qed

\section{The hypothesis: what it asserts and what it does not}\label{sec:hyp}

Hypothesis $\mathrm{H}(\vartheta)$ asks for a lower bound, with no loss against the
Buchstab value, for the $\N^{\vartheta}$-rough integers in an interval of length
$\asymp x^{1/2}$ near $x$. This section records what such a bound would and would not
require. First, the Type~I remainders of the sieve in question, against arbitrary
bounded weights, are controlled unconditionally at level $\N^{1-\varepsilon}$
(Lemma~\ref{lem:remainder}). Second, the linear sieve cannot reach the bound at any
level of distribution below $x$ (Proposition~\ref{prop:parity}), and by Selberg's
examples no argument from the sieve axioms alone can do better
(Remark~\ref{rem:twist}). Together these locate the hypothesis at the parity barrier.

Set
\[
B:=\exp\!\bigl(\sqrt{\log \N}\bigr),\qquad
S(B):=\#\{k\in\A:\ \pmin(k)>B\}.
\]

\begin{lemma}\label{lem:fund}
$S(B)=\bigl(1+O(e^{-c\sqrt{\log\N}})\bigr)\,(2n+1)\,V(B)$ for some absolute $c>0$, where
$V(B)=\prod_{p<B}(1-1/p)$. In particular $S(B)\asymp n/\sqrt{\log \N}$.
\end{lemma}

\begin{proof}
Sieve $\A$ by the primes below $B$ with the fundamental lemma at level
$D_0=\N^{\delta}$ for a small fixed $\delta>0$
\cite[Thm.~6.9]{FriedlanderIwaniec}. Since
$s_0=\log D_0/\log B=\delta\sqrt{\log \N}\to\infty$, the sieve factor is
$1+O(e^{-s_0})$, while the remainder is at most
$\sum_{d\le D_0}1\le \N^{\delta}$, because each individual remainder
$\bigl||\A_d|-(2n{+}1)/d\bigr|$ is at most $1$; and $\N^{\delta}\ll nV(B)e^{-s_0}$.
\end{proof}

For squarefree $m$ with $\pmin(m)>B$ define the discrepancy
\begin{equation}\label{eq:delta-def}
\Delta(m):=\#\Bigl\{k\in\A:\ k\equiv 0 \ (\mathrm{mod}\ m),\ \pmin(k/m)>B\Bigr\}
\;-\;\frac{S(B)}{m}.
\end{equation}
These are the remainders of the sieve of the $B$-rough elements of $\A$ by the primes
in $[B,z)$, whose moduli are the squarefree $m$ composed of such primes.

\begin{lemma}[Unconditional remainder control]\label{lem:remainder}
Fix $\vartheta\in(0,1)$. There is $c=c(\vartheta)>0$ such that
\[
\sum_{\substack{m\le \N^{\vartheta}\\ \mu^2(m)=1,\ \pmin(m)>B}} |\Delta(m)|
\ \ll\ n\,e^{-c\sqrt{\log \N}}.
\]
Consequently, for every choice of coefficients $(c_m)$ with $|c_m|\le 1$ supported on
such $m$,
\[
\Bigl|\sum_m c_m\Delta(m)\Bigr|\ \ll\ n\,e^{-c\sqrt{\log\N}}=o\bigl(u/\log\N\bigr),
\]
so every Type~I remainder of this sieve aggregated by bounded weights of level
$\N^{\vartheta}$ --- upper or lower beta-sieve weights, well-factorable weights, or any
bounded weights whatsoever --- is admissible.
\end{lemma}

\begin{proof}
Fix $\varepsilon\in(0,1/2)$. For $m$ in the sum put $I_m:=\{k'\in\mathbb{Z}:\ mk'\in\J\}$,
an interval of length $(2n+1)/m+O(1)\ge \tfrac12\N^{1-\vartheta}$. Writing $k=mk'$, the count
in \eqref{eq:delta-def} is $S(I_m,B):=\#\{k'\in I_m:\ \pmin(k')>B\}$. Sieve $I_m$ by
the primes below $B$ at level $D_m:=|I_m|^{1-\varepsilon}$: the individual remainders
are at most $1$ in absolute value, so their sum is at most $D_m$, and
$s_m:=\log D_m/\log B\ge (1-\varepsilon)(1-\vartheta)\sqrt{\log\N}\,(1+o(1))\to\infty$
uniformly in $m$. The fundamental lemma \cite[Thm.~6.9]{FriedlanderIwaniec} gives
\[
S(I_m,B)=|I_m|\,V(B)\bigl(1+O(e^{-s_m})\bigr)+O\bigl(|I_m|^{1-\varepsilon}\bigr).
\]
By Lemma~\ref{lem:fund}, $S(B)/m=(2n+1)V(B)\bigl(1+O(e^{-c\sqrt{\log\N}})\bigr)/m$, and
$|I_m|=(2n+1)/m+O(1)$. Hence
\[
\Delta(m)\ \ll\ \frac{n}{m}\,V(B)\,e^{-c'\sqrt{\log \N}}+\Bigl(\frac nm\Bigr)^{1-\varepsilon}+V(B).
\]
Summing over $m\le\N^{\vartheta}$, the three terms contribute
$\ll nV(B)L\,e^{-c'\sqrt{\log\N}}$, $\ll n^{1-\varepsilon}\N^{\varepsilon\vartheta}
=\N^{1-\varepsilon(1-\vartheta)+o(1)}$ and $\ll \N^{\vartheta}$ respectively, each
$\ll n\,e^{-c\sqrt{\log\N}}$. The consequence is the triangle inequality, and
$u/\log\N=n/L^2$.
\end{proof}

Lemma~\ref{lem:remainder} is the analogue, for $\J$, of the exact truncated-Legendre
identification in \cite[\S\S6--7]{RossShells}: on the Type~I side of the problem there
is no distributional obstacle at all. The obstacle is the following.

\begin{proposition}[$\mathrm{H}(\vartheta)$ lies beyond the sieve axioms]\label{prop:parity}
Let $\vartheta\in(2/3,1)$, $z=\N^{\vartheta}$, and
$t(\vartheta):=\log\frac{2-\vartheta}{\vartheta}$, so that $m(\vartheta)=1+t(\vartheta)$.
\begin{enumerate}
\item[(i)] The linear sieve lower bound for $S(\A,z)$ at the level $\N^{1-\varepsilon}$
of Lemma~\ref{lem:remainder} is trivial: $s=(1-\varepsilon)/\vartheta<2$ and $f(s)=0$.
\item[(ii)] Suppose that remainder control were available at some level $D=\N^{\kappa}$
with $\kappa<2$, that is, $D<x$. Then the linear sieve lower bound is at most
\[
\bigl(f(\kappa/\vartheta)+o(1)\bigr)(2n+1)V(z)\ \le\ \bigl(2t(\vartheta)+o(1)\bigr)\,u .
\]
Since $2t(\vartheta)<m(\vartheta)$ exactly when $t(\vartheta)<1$, i.e.\ for
$\vartheta>2/(1+e)=0.53788\ldots$, Hypothesis $\mathrm{H}(\vartheta)$ exceeds every
linear-sieve lower bound by $\bigl(1-t(\vartheta)-o(1)\bigr)u$, at every level of
distribution, for every $\vartheta\in(2/(1+e),1)$. At $\vartheta=\thstar$ the sieve
bound is $0.243\,u$ against the hypothesis's $1.122\,u$.
\end{enumerate}
\end{proposition}

\begin{proof}
The linear sieve lower bound reads
$S(\A,z)\ge (f(s)-o(1))(2n+1)V(z)-\sum_{d\le D}|r_d|$ with $s=\log D/\log z$
\cite{FriedlanderIwaniec}; this gives (i). For (ii), $f$ is increasing, so the bound
at level $\N^{\kappa}$ is at most the bound at $s=2/\vartheta\in(2,3)$, where
$f(2/\vartheta)=e^{\gamma_0}\vartheta\,t(\vartheta)$, while
$(2n+1)V(z)=2n\,e^{-\gamma_0}(1+o(1))/(\vartheta L)$. The product is
$2t(\vartheta)\,u\,(1+o(1))$. The numerical values are $t(\thstar)=0.12154\ldots$.
\end{proof}

\begin{remark}[Parity examples]\label{rem:twist}
Selberg's examples show that the linear-sieve bounds are best possible under the sieve axioms: nonnegative prime-free sequences can satisfy the same Type~I data and attain the lower bound $f(s)$ \cite{Selberg,FriedlanderIwaniec}. Thus Proposition~\ref{prop:parity} reflects the parity barrier rather than a weakness of the chosen sieve. A local Liouville version for the present intervals is discussed in \cite[\S\S12--14]{RossShells}.
\end{remark}

\begin{remark}[What $\mathrm{H}(\vartheta)$ requires, and its strength]\label{rem:requires}
(a) What passes the parity barrier in the known cases is bilinear, Type~II,
information about the sequence itself: cancellation, schematically, in
\[
\sum_{m\sim M}\sum_{k\sim K}\alpha_m\beta_k\bigl(\mathbf{1}_{\J}(mk)-\text{expected}\bigr),
\qquad MK\asymp x,
\]
for general bounded coefficients on \emph{both} variables and some range of $M,K$,
in the sense of the asymptotic sieve of Friedlander and Iwaniec \cite{FI98} and of
Harman's method \cite{Harman,FordMaynard}. A combinatorial decomposition of the kind
used in Harman's sieve organizes such information; it does not replace it. For
windows of length $x^{1/2}$ no Type~II estimate is known: the cited
Dirichlet-polynomial results reach primes in intervals of length
$x^{17/30+\varepsilon}$ \cite{GuthMaynard}. Absolute remainders cannot supply it: for a window
of length $Y$ and $D\ge 2Y$ one has $\sum_{D<d\le 2D}\bigl||\A_d|-Y/d\bigr|\gg Y$
\cite[\S14]{RossShells}, so the level $\N^{1-\varepsilon}$ of
Lemma~\ref{lem:remainder} is the most that Type~I information can give, and by (i) it
gives nothing.

(b) $\mathrm{H}(\vartheta)$ is stronger than prime existence: by \eqref{eq:buchstab-intro} it is a prime-number-theorem-strength lower bound for primes and the semiprime tail together. The point of the reduction is that the semiprime contribution is now controlled unconditionally.
\end{remark}

\section{The race and the threshold}\label{sec:race}

\begin{proof}[Proof of Corollary~B]
By Lemma~\ref{lem:buchstab} with $\theta=\vartheta$, Hypothesis
$\mathrm{H}(\vartheta)$, and Theorem~A,
\[
\pi(\J)= S(\A,\N^{\vartheta})-T(\vartheta)+O(1)
\ \ge\ \bigl(m(\vartheta)-C(\vartheta)-o(1)\bigr)u .
\]
It remains to note $c(\vartheta)=m(\vartheta)-C(\vartheta)>0$ for
$\vartheta>\thstar$: both sides of \eqref{eq:threshold} are continuous, the left side
is decreasing and the right side strictly decreasing to $0$ as
$\vartheta\to1^-$ faster (indeed $C(\vartheta)\to0$ while $m\to1$), while as
$\vartheta\to(4/5)^+$ the right side diverges; a single crossing follows from
monotonicity of $C-m$ (its derivative is negative on $(4/5,1)$, as
$C'(\vartheta)=-\tfrac{16}{(2-\vartheta)(5\vartheta-4)}<m'(\vartheta)
=-\tfrac{2}{\vartheta(2-\vartheta)}$ there, both computed directly). Numerically
$\thstar=0.939305\ldots$
\end{proof}

\section{Discussion}\label{sec:discussion}

\subsection{General intervals}\label{sec:general}
Sections~\ref{sec:tail-setup}--\ref{sec:expsums} use only $x\asymp\N^2$, $y\asymp\N$
and the shape of the tail, never the square center. Theorem~A therefore holds, with
the same proof, for any interval $[x-y,\,x+y]$ with $y\asymp x^{1/2}$: writing
\[
T(\theta):=\#\{m\in[x-y,x+y]:\ m=pq,\ x^{\theta/2}<q\le x^{1/2}<p\},
\qquad u:=\frac{2y}{\log x},
\]
so that $u$ is the prime number theorem's prediction for the number of primes in the
interval (for $\J$ one has $y=n$ and $u=n/\log\N$ as before), one has
$T(\theta)\le\bigl(C(\theta)+o(1)\bigr)\,u$ for every fixed $\theta\in(4/5,1)$,
against the Brun--Titchmarsh value $(4+o(1))\,u$ of
Proposition~\ref{prop:classical}. The square center enters only through
Lemma~\ref{lem:structural}, which makes the Buchstab identity
\eqref{eq:buchstab-intro} exact for $\J$ with no boundary terms.

The square center is needed only for the exact prime-detecting identity; the tail estimate itself applies to general square-root intervals. Related square-centered obstructions are studied in \cite{RossObstruction,RossRays,RossShells}. Thus the remaining input is genuinely Type~II information for intervals of length $x^{1/2}$: Theorem~A removes the semiprime-tail obstruction but does not supply that parity-breaking input.

\appendix

\section{Exact certification of the level profile}\label{app:cert}

The optimization behind Lemma~\ref{lem:coverage} and Proposition~\ref{prop:optimal} is piecewise linear and was checked exactly. For each fixed $\gamma$, admissibility of the Robert--Sargos groupings, the additional \cite[Thm.~6]{FI} groupings, and the elementary exponent-pair bounds becomes a finite system of linear inequalities in $(\mu,\eta,\beta)$. Exact rational vertex enumeration gives
\[
\begin{array}{c|ccccc}
\gamma & 1 & 21/20 & 11/10 & 23/20 & 6/5\\ \hline
\zeta(\gamma) & 1/4 & 3/16 & 1/8 & 1/16 & 0
\end{array}
\qquad\Bigl(=\tfrac{6-5\gamma}{4}\Bigr).
\]
In every case the minimizing witness is the balanced type-II block at top frequency, $(\mu,\eta)=(\gamma/2,\Sigma-1)$, confirming the facet $4\beta+5\gamma=6$. The computation also shows that adding the extra Fouvry--Iwaniec groupings and the exponent pair $(\tfrac16,\tfrac23)$ does not improve the optimum. Code and the full constraint inventory are archived at
\url{https://github.com/michaelmross/Legendre/tree/main/Jn-tail}; the table is produced by \texttt{fm\_exact.py}.

\subsection*{Acknowledgments}
The author thanks Claude Fable 5.1 (Anthropic) and GPT-5.6 Sol (OpenAI) for collaboration in analysis and exposition.

\begin{thebibliography}{99}

\bibitem{FordMaynard} K.~Ford and J.~Maynard,
\emph{On the theory of prime producing sieves},
arXiv:2407.14368 (2024).

\bibitem{FI} E.~Fouvry and H.~Iwaniec,
\emph{Exponential sums with monomials},
J. Number Theory \textbf{33} (1989), 311--333.

\bibitem{FI98} J.~Friedlander and H.~Iwaniec,
\emph{Asymptotic sieve for primes},
Ann. of Math. (2) \textbf{148} (1998), 1041--1065.

\bibitem{FriedlanderIwaniec} J.~Friedlander and H.~Iwaniec,
\emph{Opera de Cribro},
Amer. Math. Soc. Colloq. Publ. \textbf{57}, AMS, 2010.

\bibitem{GK} S.~W. Graham and G.~Kolesnik,
\emph{Van der Corput's Method of Exponential Sums},
London Math. Soc. Lecture Note Ser. \textbf{126}, Cambridge Univ. Press, 1991.

\bibitem{GuthMaynard} L.~Guth and J.~Maynard,
\emph{New large value estimates for Dirichlet polynomials},
Ann. of Math. (2) \textbf{203} (2026), 623--675.

\bibitem{Harman} G.~Harman,
\emph{Prime-Detecting Sieves},
London Math. Soc. Monographs \textbf{33}, Princeton Univ. Press, 2007.

\bibitem{Iwaniec80} H.~Iwaniec,
\emph{A new form of the error term in the linear sieve},
Acta Arith. \textbf{37} (1980), 307--320.

\bibitem{RS} O.~Robert and P.~Sargos,
\emph{Three-dimensional exponential sums with monomials},
J. reine angew. Math. \textbf{591} (2006), 1--20.

\bibitem{RossRays} M.~M. Ross,
\emph{Factor Rays and the Self-Conjugate Parabola: Deterministic Coverage Geometry in
Square Intervals},
Zenodo (2026), \doi{10.5281/zenodo.20016398}.

\bibitem{RossShells} M.~M. Ross,
\emph{Multiplication Geometry in Square Shells: Finite Coverage, Truncated Legendre
Sums, and the Parity Barrier},
Zenodo (2026), \doi{10.5281/zenodo.22553216}.

\bibitem{RossObstruction} M.~M. Ross,
\emph{Primes in Square Intervals: Obstructions to the Square-Centered Dispersion Route},
Zenodo (2026), \doi{10.5281/zenodo.21248789}.

\bibitem{Selberg} A.~Selberg,
\emph{Sieve methods},
in: 1969 Number Theory Institute, Proc. Sympos. Pure Math. \textbf{XX},
Amer. Math. Soc., Providence, 1971, pp.~311--351.

\bibitem{Tenenbaum} G.~Tenenbaum,
\emph{Introduction to Analytic and Probabilistic Number Theory}, 3rd ed.,
Grad. Stud. Math. \textbf{163}, AMS, 2015.

\bibitem{Vaughan} R.~C. Vaughan,
\emph{Sommes trigonom\'etriques sur les nombres premiers},
C. R. Acad. Sci. Paris S\'er. A \textbf{285} (1977), 981--983.

\end{thebibliography}

\end{document}